\documentclass[
    12pt,                              % Schriftgröße
    fleqn,                             % Gleichungen linksbündig
    twoside,
    headsepline,                       % einzeichnen einer Linie unter die Kopfzeile
    listof=totoc,
    toc=chapterentrywithdots,
    bibliography=totocnumbered,        % Nummerierung des Literaturverzeichnis
    parskip=half,                      % Absatzabstände erzeugen
]{scrreprt}

\usepackage[                            % Seitenränder
    left = 25mm,
    right = 25mm,
    top = 30mm,
    bottom = 25mm,
    bindingoffset=15mm,
]{geometry}

\usepackage[pdfpagelabels,
    pdfstartview = FitH,
    bookmarksopen = false,
    bookmarksnumbered = true,
    linkcolor = black,
    plainpages = false,
    hypertexnames = false,
    citecolor = black
]{hyperref}

\newcommand*{\captionsource}[2]{%
  \caption[{#1}]{%
    #1%
    \\\hspace{\linewidth}%
    Quelle: #2%
  }%
}


\usepackage[T1]{fontenc}
\usepackage[utf8]{inputenc}                 % UTF8 (Windowsspezifisch), Für Umlaute wie ä,ö,ü,$,...
\usepackage[ngerman]{babel}                 % Deutsch
\usepackage{lmodern}                        % Schriftart 
\usepackage{graphicx, subfig}
\usepackage{float}
\usepackage{scrhack}                        % Warnung wegen floats unterdrücken
\usepackage{pdfpages}
\graphicspath{{img/}}
\setkomafont{disposition}{\bfseries}        % Dei Schriftart für alles setzen (Überschriften etc.)
\usepackage[onehalfspacing]{setspace}       % Zeilenabstand 1,5cm
\BeforeStartingTOC[toc]{\singlespacing}     % "Korrekter" Zeilenabstand im Inhaltsverzeichnis

\usepackage{scrlayer-scrpage}               % Um Kopf- und Fußzeile ändern zu können
\clearpairofpagestyles                      % Löscht Standardeinstellungen

\ihead{\normalfont{\headmark}}                  % Header Inside
\automark[chapter]{chapter}                     % \automark[rechte Seite]{linke Seite}
\chead{}                                        % Header Center
\ohead{\pagemark}                               % Header Outside
\renewcommand*\chapterpagestyle{scrheadings}    % Header auch für \chapter generieren
    % Wie oben nur für Footer
\ifoot{\normalfont\footnotesize{Bastian Huber}}
\cfoot{}
\ofoot{\normalfont\footnotesize{\today}}

\usepackage{lipsum}                         % Blindtext mit \lipsum generieren
                    
\addtokomafont{chapter}{\large}             % Chapter, Section etc. Schriftgröße
\addtokomafont{section}{\normalfont\large}
\addtokomafont{subsection}{\normalfont\normalsize}

%#################################################################
% Begin des Document
% ----------------------------------------------------------------
\begin{document}
    \pagenumbering{Roman}
    %#################################################################
    % Titelseite
    % ----------------------------------------------------------------
    \newgeometry{           % Custom Seitenlayout für die Titelseite (kein Bund)
    left = 25mm,
    right = 25mm,
    top = 30mm,
    bottom = 25mm,
    bindingoffset=0mm,
}
\thispagestyle{empty}
\begin{titlepage}
    \hfill \includegraphics[scale=0.8]{1_Titelseite/LogoHFU.pdf}\\
    \begin{tabular}{p{\textwidth}}
        \begin{center}
            \LARGE{[Art der Arbeit]}
        \end{center}
        
        \begin{center}
            \LARGE{[Studiengang]}
        \end{center}
        \\
        \begin{center}
            \textbf{\Huge{[Titel]}}
        \end{center}

        \begin{center}
            \LARGE{[Untertitel]}
        \end{center}
        
        \begin{minipage}{\textwidth}
            \vskip 8cm
            {\normalsize}
            \begin{center}
                \begin{tabular}{ll}
                    Referent: & [Name]   \tabularnewline
                    Korreferent:& [Name]   \tabularnewline
                    Vorgelegt am:& [Datum]   \tabularnewline
                    Vorgelegt von:& [Name] \tabularnewline
                        & [Matrikelnummer]   \tabularnewline
                        & [Straße]   \tabularnewline
                        & [PLZ Ortschaft]   \tabularnewline
                        & [Email]    \tabularnewline
                \end{tabular}
            \end{center}
            {\normalsize}
        \end{minipage}
    \end{tabular}
\end{titlepage}
\restoregeometry        % Seitenlayout wird wieder zurückgesetzt
  
    \cleardoublepage
    
    %#################################################################
    % Abstract und Vorwort
    % ----------------------------------------------------------------
    \addchap{Abstract}
\lipsum[1]
    \cleardoublepage

    %#################################################################
    % ToC
    % ----------------------------------------------------------------
    \tableofcontents
    \cleardoublepage 

    \listoffigures
    \cleardoublepage 

    %\listoftables
    %\cleardoublepage                  % zum Beenden der roman Seitennummerierung! \clearpage

    \pagenumbering{arabic}            % Seitennummerierung Hauptteil
    %#################################################################
    % Inhalt
    % ----------------------------------------------------------------                      
    \chapter{Einleitung}
\lipsum[1]
    \cleardoublepage
    \chapter{Kapitel 1}
\lipsum[2]
\section{Unterkapitel 1.1}
\lipsum[3]\cite{Tran_2008}
\section{Unterkapitel 1.2}
\lipsum[4]
\chapter{Kapitel 2}
\lipsum[5]
\section{Unterkapitel 2.1}
\lipsum[6]
\includegraphics*{images/HFU.jpg}
\chapter{Kapitel 3}
\lipsum[7]
    \cleardoublepage
    \chapter{Schluss}
\lipsum[8]

    \cleardoublepage

    %#################################################################
    % Anhaenge
    % ----------------------------------------------------------------                      
    \addchap{Eidesstattliche Erklärung}
\label{erklaerung}

Hiermit versichere ich, die vorliegende Arbeit selbstständig und nur unter Verwendung der von mir angegebenen Quellen und Hilfsmittel verfasst zu haben. Sowohl inhaltlich als auch wörtlich entnommene Inhalte wurden als solche kenntlich gemacht. Die Arbeit hat in dieser oder vergleichbarer Form noch keinem anderem Prüfungsgremium vorgelegen. \\
\\[1.5cm]
Datum:	\hrulefill\enspace Unterschrift: \hrulefill
\\[3.5cm]

    %\raggedright                       % Blocksatz ausstellen
    %\bibliographystyle{unsrt}           % Literatur nummeriert darstellen
    \bibliographystyle{plain}
    \bibliography{Literaturstellen}     % .bib einbinden

\end{document}
